\chapter{Urinary Tract}

\label{ch:ch50_urinary_tract}

\section{What is the urinary tract?}
The urinary tract consists of the kidneys, ureters, bladder, and urethra. Its role is to produce, transport, store, and expel urine. Imaging of the urinary tract is among the most commonly performed radiological investigations, with kidney stones, haematuria, and obstruction representing the bulk of referrals.

\section{When is urinary tract imaging ordered?}
\begin{itemize}
    \item \textbf{Suspected kidney stones (renal colic):} Sudden, severe flank pain radiating to the groin, often with blood in the urine.
    \item \textbf{Haematuria:} Blood in the urine without pain may indicate a bladder or ureteric tumor.
    \item \textbf{Urinary obstruction:} Inability to pass urine, flank pain, or rising creatinine.
    \item \textbf{Recurrent urinary tract infections:} Structural abnormality or stone disease may be responsible.
    \item \textbf{Trauma:} Bladder or ureteric injury after pelvic fractures or penetrating trauma.
    \item \textbf{Known malignancy:} Staging and surveillance of bladder, ureteric, or renal pelvis tumors.
\end{itemize}

\section{What happens during the scan?}

\subsection{CT KUB (non-contrast CT of kidneys, ureters, bladder)}
This is the investigation of choice for suspected renal colic.
\begin{itemize}
    \item You lie on a table that passes through the CT scanner.
    \item No contrast injection is needed.
    \item The scan takes 2--3 minutes.
    \item It detects stones with $>$95\% sensitivity and specificity.
    \item It also identifies alternative diagnoses (appendicitis, ovarian pathology, aortic pathology) in up to 30\% of cases of suspected colic.
\end{itemize}

\subsection{CT urography}
This is a contrast-enhanced study that evaluates the entire urinary tract for masses, haematuria, and structural abnormalities. It involves pre-contrast, nephrographic, and excretory phases. You will receive an IV contrast injection and may be asked to delay imaging while the contrast is excreted into the collecting system.

\subsection{Ultrasound}
Ultrasound is the first-line investigation for:
\begin{itemize}
    \item Assessment of hydronephrosis
    \item Bladder volume and post-void residual
    \item Renal masses in children
    \item Pregnancy-related hydronephrosis (no radiation)
\end{itemize}

\section{Urolithiasis: imaging findings}

\begin{table}[h]
\centering
\caption{Imaging modalities for urinary stones}
\label{tab:stone_modalities}
\begin{tabular}{@{}p{0.15\textwidth}p{0.12\textwidth}p{0.32\textwidth}p{0.25\textwidth}@{}}
\hline
\textbf{Modality} & \textbf{Sensitivity} & \textbf{Advantages} & \textbf{Limitations} \\
\hline
CT KUB & $>$95\% & Gold standard, identifies stone size/location & Radiation exposure \\
IVU & $>$90\% & Shows function + anatomy & Largely replaced by CT \\
KUB X-ray & 44--77\% & Low dose, quick & Misses uric acid and small stones \\
Ultrasound & 45--75\% & No radiation, portable, good for hydronephrosis & Misses ureteric stones \\
\hline
\end{tabular}
\end{table}

\section{Stone composition clues on CT}
Stone density on CT (measured in Hounsfield units) provides a clue to composition, which may guide management:

\begin{table}[h]
\centering
\caption{CT density and stone composition}
\label{tab:stone_density}
\begin{tabular}{@{}p{0.24\textwidth}p{0.20\textwidth}p{0.48\textwidth}@{}}
\hline
\textbf{Stone type} & \textbf{CT density (HU)} & \textbf{Clinical note} \\
\hline
Calcium oxalate & 1200--1800 & Most common stone type \\
Calcium phosphate & 1200--1600 & Alkaline urine \\
Cystine & 600--1100 & Rare; genetic; hexagonal crystals on urine microscopy \\
Struvite (infection) & 600--900 & Associated with urease-producing bacteria \\
Uric acid & 200--500 & Radiolucent on X-ray; dissolves with alkalinisation \\
\hline
\end{tabular}
\end{table}

\begin{pearlbox}[Uric acid stones and CT]
Uric acid stones are the only common stone type that can be dissolved medically (with potassium citrate to alkalinise urine). They are often invisible on plain X-ray but visible on CT as moderate-density filling defects. Look for the ``rim sign''---a rim of higher-density calcium phosphate surrounding a lower-density uric acid stone.
\end{pearlbox}

\section{When to order CT vs ultrasound for suspected stones}
\begin{itemize}
    \item \textbf{CT KUB is preferred} for: acute renal colic in adults, suspected complicated stone disease, equivocal ultrasound, and pre-intervention planning.
    \item \textbf{Ultrasound is preferred} for: pregnant patients, children, initial assessment of hydronephrosis, follow-up of known stones, and screening.
    \item \textbf{Consider low-dose CT} for: young patients, follow-up imaging, and screening for recurrence (reduces dose by 50--80\%).
\end{itemize}

\section{Ureteric jets on ultrasound}
Color Doppler ultrasound can demonstrate ureteric jets---intermittent bursts of urine entering the bladder from the ureteric orifices. The presence of a ureteric jet does not exclude partial obstruction, but the \textbf{absence of a ureteric jet on the affected side} in the setting of hydronephrosis strongly suggests complete or near-complete obstruction. Asymmetric jet frequency or amplitude may indicate partial obstruction.

\begin{pearlbox}[Twinkling artefact]
A color Doppler artefact called ``twinkling'' may appear behind ureteric stones as a rapidly alternating color signal. This can help confirm the presence of a small stone that might otherwise be difficult to identify, particularly on ultrasound.
\end{pearlbox}

\section{Ureteric obstruction: causes and imaging}
Obstruction at different levels of the ureter suggests different causes:

\begin{table}[h]
\centering
\caption{Causes of ureteric obstruction by location}
\label{tab:obstruction_causes}
\begin{tabular}{@{}p{0.17\textwidth}p{0.36\textwidth}p{0.33\textwidth}@{}}
\hline
\textbf{Location} & \textbf{Common causes} & \textbf{Key imaging finding} \\
\hline
Mid-ureter & Stones, retroperitoneal fibrosis, lymphadenopathy & Dilation above, collapsed below \\
PUJ (proximal) & Congenital PUJ obstruction, stones, crossing vessel & Pelvicalyceal dilation, no ureteric dilation \\
UVJ (distal) & Stones (most common), bladder tumor, endometriosis & Ureteric wall thickening, filling defect \\
\hline
\end{tabular}
\end{table}

\section{Bladder masses}
Bladder masses are typically identified on CT urography or present as haematuria. The most common bladder malignancy is urothelial (transitional cell) carcinoma ($>$90\% of bladder cancers). On imaging:
\begin{itemize}
    \item \textbf{Papillary urothelial carcinoma:} A polypoid mass projecting into the bladder lumen.
    \item \textbf{Sessile/infiltrating carcinoma:} Irregular wall thickening with possible invasion into perivesical fat.
    \item \textbf{Staging} requires assessment of depth of invasion (T1--T4), lymph nodes, and distant metastases. MRI is superior for local staging (T2-weighted imaging shows muscle invasion).
\end{itemize}

\section{Urinary diversion imaging}
Patients who have undergone radical cystectomy require imaging surveillance of their urinary diversion:
\begin{itemize}
    \item \textbf{Ileal conduit:} A segment of small bowel is used to channel urine to a stoma. Imaging evaluates for ureteroconduit anastomotic strictures, conduit strictures, stones within the conduit, and recurrence.
    \item \textbf{Neobladder:} A bowel segment is configured into a bladder-like reservoir. CT evaluates capacity, ureteric anastomosis, and complications.
    \item \textbf{Cutaneous ureterostomy:} Direct ureter-to-skin stoma; evaluate for stenosis and obstruction.
\end{itemize}

\section{Reporting checklist}
A structured urinary tract report should include:
\begin{enumerate}
    \item Kidneys: size, hydronephrosis grade, parenchymal changes
    \item Ureters: course, dilatation, transition point, filling defects
    \item Stones: number, size (mm), location (renal, mid-ureter, UVJ), degree of obstruction
    \item Bladder: wall thickness, masses, diverticula, calculi, post-void residual
    \item Prostate (in males): size, intravesical protrusion
    \item Uterus/ovaries (in females): relevant adnexal pathology
    \item Free fluid
    \item Alternative diagnoses
\end{enumerate}

\clearpage
